% Owner: ________
% Scope: GUI agents, mobile environments, hybrid GUI + native tool use
\section{Related Work}
\label{sec:related}

\paragraph{GUI Agents.}
Early GUI agents operated on structured representations such as HTML, accessibility trees, or view hierarchies. These signals can be long, incomplete, or unavailable on some platforms, motivating purely visual agents that act directly from screenshots~\citep{cheng2024seeclick,hong2024cogagent}. Subsequent work improved GUI grounding and cross-platform action modeling~\citep{wu2025osatlas}, while mobile-specific systems demonstrated that multimodal models can complete multi-step phone tasks through tapping, typing, and swiping~\citep{yang2023appagent,wang2024mobileagent}. More recent foundation agents have further improved mobile task performance through large-scale training and online reinforcement learning~\citep{liu2024autoglm,qin2025uitars}. Despite these advances, the dominant interaction paradigm remains GUI-based: agents perceive pixels, optionally augmented with UI structure, and emit low-level actions. 

\paragraph{Mobile Environments.}
Offline corpora such as AITW and AndroidControl provide large-scale human demonstrations for mobile interaction~\citep{rawles2023aitw,li2024androidcontrol}. Interactive environments have subsequently enabled agents to execute and evaluate tasks directly on Android devices and emulators. AndroidWorld provides a widely used open environment with 116 programmatic tasks across 20 Android applications, featuring parameterized instructions and state-based rewards~\citep{rawles2025androidworld}. AndroidLab and SPA-Bench further expand the range of tasks and application settings, while MobileWorld pushes toward more challenging long-horizon and cross-application workflows and introduces MCP-augmented interaction~\citep{xu2025androidlab,chen2024spabench,kong2025mobileworld}. Complementing these device- and emulator-based environments, MobileGym takes a simulation-based approach by recreating application experiences in a browser-hosted sandbox and mocking external dependencies, enabling commercial-style workflows to be reset, programmatically verified, and rolled out in parallel~\citep{wu2026mobilegym}. Together, these efforts have progressively expanded the scale, diversity, and executability of mobile-agent environments, providing increasingly strong foundations for both training and evaluation.

\paragraph{Hybrid GUI and App-Native Tool Use.}
Beyond low-level GUI actions, recent work has explored richer interfaces for software interaction. On desktop, UFO and UFO$^2$ combine UI automation with application-level APIs, while OS-Copilot exposes heterogeneous capabilities such as bash, Python, and APIs alongside mouse and keyboard interaction~\citep{zhang2024ufo,zhang2025ufo2,wu2024oscopilot}. The Model Context Protocol (MCP) provides a standardized protocol for exposing and invoking tools across applications~\citep{anthropic2024mcp}. In the mobile domain, MobileWorld incorporates MCP-based tool interaction into its evaluation tasks~\citep{kong2025mobileworld}. These approaches extend the action space beyond direct GUI manipulation by allowing agents to invoke structured application or system capabilities. A related setting is to expose selected application functionalities as app-native tools alongside the GUI, allowing GUI actions and tool calls to be interleaved within the same task while retaining the GUI as a complete interaction interface.

